% Šablona prezentace s aplikací nového vizuálního stylu TUL
% Upraveno pro semestrální práci – Sběr dat a monitorování procesů
\newcommand{\verze}{Verze 2.3}

\newcommand{\TULextlogofile}{V21TUL_pruhledne_CZ.pdf}

% NASTAVENÍ: Fakulta mechatroniky (FM), poměr stran 16:9
\documentclass[FM,aspect169]{tulpresentation} 

\usepackage{polyglossia}
\setdefaultlanguage{czech}

% HLAVIČKA DOKUMENTU – TVÉ ÚDAJE
\title{Sběr dat a monitorování procesů pomocí vestavěných systémů}
\author{David Švancar}
\institute{FM TUL}

\newcommand{\authorMail}{david.svancar@tul.cz}
\date{11. června 2026}

\begin{document} 

% 1. Úvodní snímek (Povinný dle zadání i šablony)
\TULtitleframe

% 2. Obsah prezentace / Outline (Povinný dle zadání)
\begin{frame}{Obsah prezentace} 
  \tableofcontents[sections={1-}]
\end{frame} 

% 3. Vlastní prezentace tématu (Max. 5 snímků dle zadání)

\section{Úvod a cíle}
% Snímek 1: Úvod
\begin{frame}{Úvod a cíle}
  \begin{itemize}
    \item \textbf{Průmyslová automatizace:} Rostoucí nároky na efektivní sběr a vyhodnocování procesních dat v reálném čase.
    \item \textbf{Vestavěné systémy:} Klíčový stavební kámen díky vysoké spolehlivosti a nízké spotřebě energie.
    \item \textbf{Cíl semestrální práce:} Návrh a realizace modulární monitorovací jednotky pro sběr fyzikálních veličin v náročném průmyslovém prostředí.
  \end{itemize}
\end{frame}

\section{Hardwarová architektura}
% Snímek 2: Hardware
\begin{frame}{Hardwarová architektura}
  \begin{block}{Řídicí jádro}
    32bitový mikrokontrolér s procesorovým jádrem \textbf{ARM Cortex-M4} (zajišťuje dostatečný výkon pro filtraci signálů).
  \end{block}
  \begin{itemize}
    \item \textbf{Napájecí část:} Stabilizace napětí z průmyslových $24\,\mathrm{V}$ na logických $3,3\,\mathrm{V}$.
    \item \textbf{Senzorické rozhraní:} Integrované ADC převodníky pro vyhodnocování napěťových úrovní ze senzorů teploty a tlaku.
    \item \textbf{Komunikační modul:} Průmyslová sběrnice RS-485 pro spolehlivý přenos dat do nadřazeného systému.
  \end{itemize}
\end{frame}

\section{Naměřené výsledky}
% Snímek 3: Data a tabulka
\begin{frame}{Analýza spotřeby a filtrace}
  Proudový odběr mikrokontroléru roste lineárně s frekvencí (viz tabulka):
  \begin{center}
    \begin{tabular}{|c|c|c|c|}
        \hline
        \textbf{Měření} & \textbf{Frekvence [MHz]} & \textbf{Napětí [V]} & \textbf{Odběr [mA]} \\ \hline
        1 & 8  & 3,3 & 4,2  \\ \hline
        4 & 96 & 3,3 & 35,1 \\ \hline
    \end{tabular}
  \end{center}
  \vspace{0.1cm}
  Vstupní obvody byly doplněny o analogový RC filtr s mezním kmitočtem $f_0 = 1023\,\mathrm{Hz}$ pro potlačení nízkofrekvenčního šumu.
\end{frame}

\section{Programová implementace}
% Snímek 4: Ukázka kódu
\begin{frame}[fragile]{Programová implementace v jazyce C}
  Konfigurace interního registru periferie ADC (přímé mapování do paměti):
  \begin{verbatim}
  #define ADC_CR_REG  (*((volatile uint32_t*)0x40012000))
  #define ADC_ENABLE  (1 << 0)

  void init_adc_peripheral(void) {
      // Zapnuti periferie ADC pomoci prislusneho bitu
      ADC_CR_REG |= ADC_ENABLE;
      
      /* Cekani na stabilizaci referencniho napeti */
      for(volatile int i = 0; i < 1000; i++);
  }
  \end{verbatim}
\end{frame}

\section{Závěr}
% 4. Závěr (Povinný dle zadání)
\begin{frame}{Závěr}
  \begin{itemize}
    \item \textbf{Ověření funkčnosti:} Navržená monitorovací stanice úspěšně integruje senzorická pole a bezpečně přenáší data.
    \item \textbf{Optimalizace:} Analýza spotřeby ukázala nutnost dynamického řízení frekvence pro energeticky úsporné režimy.
    \item \textbf{Splnění cílů:} Ukázkový dokument i prezentace byly kompletně zpracovány v systému \LaTeX\ při striktním dodržení pravidel české technické typografie TUL.
  \end{itemize}
\end{frame}

% 5. Poděkování za pozornost (Povinné dle zadání)
% Makro \TULendframe automaticky vygeneruje oficiální závěrečný slide TUL s poděkováním
\section*{Konec}
\TULendframe 

\end{document}